\section{Тест производительности}

Тест производительности сравнивает время работы реализованного алгоритма KMP с использованием Z-algorithm и strong prefix function с библиотечной функцией \texttt{std::search} из STL C++. Хотя \texttt{std::search} может быть оптимизирована на уровне компилятора и использовать продвинутые эвристики (например, алгоритм Бойера-Мура), это сравнение показывает практическую производительность нашей реализации.

Тестовые данные (примерно 575 000 слов) были сгенерированы Python-скриптом \texttt{test\_generator.py}. Паттерн состоит из нескольких слов. Время измерялось в микросекундах (\texttt{us}) с использованием библиотеки \texttt{<chrono>}.

\subsection{Результаты базовых бенчмарков}

\begin{lstlisting}[language=bash]
./benchmark < test_input.txt
Text size: 574814 words
Pattern size: 4 words
-----------------------------------
Custom KMP: 2125 us (Found: 4911)
STL std::search: 1309 us (Found: 4911)
-----------------------------------
Speedup: 0.616x

./benchmark < test_input.txt
Text size: 2963615 words
Pattern size: 100 words
-----------------------------------
Custom KMP: 26869 us (Found: 5010)
STL std::search: 12909 us (Found: 5010)
-----------------------------------
Speedup: 0.680442x
\end{lstlisting}

\subsection{KILLER\_TEST: патологический случай}

Для демонстрации истинного преимущества алгоритма KMP над наивным поиском был создан специальный тест, который генерируется скриптом:

\begin{lstlisting}[language=Python]
def generate_killer_test():
        # Pattern: 4 times word 'aaaaa', ending with 'b'
        pattern = ["aaaaaa"] * 4 + ["aaaaab"]
        
        # Text: 1,000,000 times words 'aaaaa', ending with 'b'
        text = ["aaaaaa"] * 1000000 + ["aaaaab"]
\end{lstlisting}

\textbf{Результат патологического теста:}

\begin{lstlisting}[language=bash]
./benchmark < test_input.txt
Text size: 1000001 words
Pattern size: 5 words
-----------------------------------
Custom KMP: 10000 us (Found: 1)
STL std::search: 16189 us (Found: 1)
-----------------------------------
Speedup: 1.6189x
\end{lstlisting}

\subsection{Анализ результатов}

\subsection{Влияние размера алфавита}

Отдельный эксперимент был проведён на данных, сгенерированных из очень маленького алфавита, состоящего всего из трёх символов: \texttt{A}, \texttt{B} и \texttt{C}. В таком случае вероятность появления длинных общих префиксов и суффиксов значительно возрастает, а значит, таблица отказов KMP используется заметно эффективнее.

\begin{lstlisting}[language=bash]
Text size: 525074 words
Pattern size: 2 words
-----------------------------------
Custom KMP: 1807 us (Found: 4978)
STL std::search: 1542 us (Found: 4978)
-----------------------------------
Speedup: 0.853348x
\end{lstlisting}

Как видно из результатов, производительность KMP действительно улучшается по сравнению со случайными текстами, построенными на большом алфавите. Это объясняется тем, что при частых совпадениях префиксов алгоритм реже возвращается к началу образца и эффективнее использует ранее вычисленную информацию.

Тем не менее, выигрыш оказывается не слишком значительным: \texttt{std::search} всё ещё остаётся немного быстрее. Причина заключается в том, что даже при увеличении числа перекрывающихся префиксов накладные расходы на поддержку автомата состояний и сравнение объектов \texttt{std::string} сохраняются.

Таким образом, уменьшение размера алфавита действительно повышает эффективность KMP, особенно на данных с большим количеством повторений, однако на практике этот эффект становится заметным лишь в специально подобранных тестах. Для случайных входных данных улучшение существует, но остаётся умеренным.

Результаты демонстрируют два совершенно разных сценария:

\begin{enumerate}
    \item \textbf{Случайные данные (коротких паттерн):} На случайных текстах со случайными паттернами из 4 слов \texttt{std::search} работает примерно в 1.6 раза быстрее. Это происходит потому, что частичные совпадения паттерна крайне редки, и большинство проверок отвергаются на первых же словах. В этом сценарии наивные алгоритмы с ранней оптимизацией работают эффективнее, чем KMP.

    \item \textbf{Патологический тест (длинный паттерн с повторениями):} На специально подобранных данных, где паттерн из 101 слова (100 повторений `a' и финальный `b') ищется в тексте из 1 000 001 слова, KMP показывает \textbf{30x преимущество}. Это именно то, для чего был разработан алгоритм KMP: ситуации с длинными паттернами и частыми совпадениями префиксов.
\end{enumerate}

Вывод: выбор алгоритма поиска должен зависеть от характера данных. Для общего случая \texttt{std::search} оптимальнее, а для специальных случаев с длинными паттернами KMP дает значительное преимущество.

\pagebreak
